\scanpage{339}\mspage{255}
% l. 1
soient fonctorielles en $j$.

% l. 2–8
On va faire une construction transfinie dans les ordinaux $\alpha$ tels que
\[
 \operatorname{Card}\alpha\leq\unsure{\pi|\kappa}
\]
d'un système inductif
\[
 (X_{\alpha}(\unsure{j|g}))_{\alpha\deleted{\unsure{j|J}}}
\]
d'objets de $\prod_{j\in\operatorname{Ob}J}\unsure{\mathcal E|\mathscr E}_j$. On va d'abord ignorer la donnée éventuelle (4.10.6.2).

% l. 9–12
On posera
\[
 \deleted{X_0(j)=X_j}\qquad X_0(j)=X(j), \tag*{(4.10.6.4)}
\]
et on va donner la prescription pour passer de $X_{\alpha}$ à $X_{\alpha+1}$. La construction \deleted{\unsure{sera telle}} \unsure{inclura} aussi, pour tte flèche $\unsure{\lambda|\alpha}$ \inserted{$:j'\unsure{\mathbin{-}|\longrightarrow}j\deleted{\unsure{'}}$} de $J$ et tout $\alpha$, la donnée de
% l. 13–16
\[
 X_{\alpha}(j')\xrightarrow{\unsure{t_{\alpha}|\tau_{\alpha}}(\lambda)}X_{\alpha+1}(j\deleted{\unsure{'}}) \tag*{(4.10.6.5)}
\]
\margin{flèche (dans $\mathcal E$) \unsure{au} dessus de $\lambda$}
et ces flèches, et les flèches de transition $X_{\alpha}(j)\overset{\unsure{\tau|T}_{\alpha}(j)}{\unsure{\mathbin{-}|\longrightarrow}} X_{\deleted{\unsure{j}}\unsure{\alpha+1|\alpha,}}(j)$, doivent
% archive sheet number: 338
